\clearpage
\begin{trapbox}
\textbf{\textcolor{CTrap}{易错一：把“外层有定义”当成“外层一直同向变化”}}

\textbf{错误想法：}只要 $g(x)$ 的结果能代入 $f$，就能直接用口诀。

\textbf{为什么错：} $f(t)=t^2$ 对所有 $t$ 都有定义，但负数段递减、正数段递增。
例如 $(x-1)^2$ 就必须在 $x=1$ 分开判断。

\textbf{正确做法：}检查两件事：\textbf{能不能代入；代进去的那些数，外层是不是始终同一个增减方向。}

\medskip
\textbf{\textcolor{CTrap}{易错二：忘了先求定义域}}

例如 $\ln(x^2-2x-3)$ 并非对所有 $x$ 都有定义。
只有 $x<-1$ 或 $x>3$ 才能计算。
\textbf{正确顺序：}先求定义域，再分段，再套同增异减。

\medskip
\textbf{\textcolor{CTrap}{易错三：把中间变量 $t$ 的区间写成 $x$ 的区间}}

外层 $f(t)=t^2$ 的分界位置是 $t=0$；
若 $t=x-1$，对应的自变量分界位置\textbf{是 $x=1$，不是 $x=0$}。
做题最后要把所有区间写回 $x$。

\medskip
\textbf{\textcolor{CTrap}{易错四：以为两层都递减，整体就递减}}

拿最简单的 $10-(10-x)=x$ 检查：两层递减，整体反而递增。
正确理解不是“负面的变化叠加”，而是\textbf{比较大小的方向被反转了两次}。

\medskip
\textbf{\textcolor{CTrap}{易错五：随意合并两段单调区间}}

在 $(-\infty,-1)$ 上递减，并不意味着和 $(3,+\infty)$ 可以合成一整段，
因为中间的函数值根本不存在。
同样，一个函数先减后增，不能把两侧直接说成“整段都增”或“整段都减”。

\medskip
\end{trapbox}
\clearpage
\begin{trapbox}
\textbf{\textcolor{CTrap}{易错六：只知道“单调”，就声称“严格单调”}}
\begin{center}\begin{tikzpicture}[font=\small,>=stealth,line join=round]
\draw[gray!40,rounded corners=3pt] (0,0) rectangle (14.2,9.5);
\node at (7.1,8.9) {整段不变与某一点切线水平，不能混为一谈};

\begin{scope}[shift={(4.0,4.4)},x=1.7cm,y=.85cm]
\draw[->,gray] (-1.7,0)--(1.9,0) node[right]{$x$};\draw[->,gray] (0,-3.2)--(0,3.8) node[above]{$y$};
\draw[blue,very thick,domain=-1.4:1.4,samples=140] plot(\x,{\x^3});\draw[violet,dashed,thick] (-1.5,3)--(1.5,3);
\draw[orange!85!black,thick] (-.4,0)--(.4,0);\fill[orange!85!black] (0,0) circle(2.2pt);
\end{scope}
\node[align=left,anchor=west] at (8.3,5){\textcolor{blue}{$y=x^3$：仍严格递增}\\[10pt]$f'(0)=0$ 只是一点水平\\[16pt]\textcolor{violet}{$y=3$：整段不变}\\[10pt]不能称为严格递增\\[14pt]严格单调要看整个区间};

\node[text width=12.8cm,align=center] at (7.1,.65) {“导数为零”与“不再递增”并不等价；常函数与孤立的水平切线也不能等同。};
\end{tikzpicture}\end{center}



例如 $g(x)=x$ 递增，$f(t)=3$ 在所有 $t$ 上都保持不变，
得到 $f(g(x))=3$。它不会随着 $x$ 的增大而严格增大。
因此：\textbf{口诀对严格增减，须有严格单调的前提；允许相等时，只能相应地判断不减或不增。}

\medskip
\textbf{\textcolor{CTrap}{易错七：导数为零就断言“停止递增”}}

例如 $x^3$ 在 $x=0$ 处的导数等于 $0$，
但 $x^3$ 在整个实数范围内仍严格递增。
\textbf{导数法是可选的快速检查，最根本的依据仍是任意两点函数值的比较。}

\medskip
\textbf{\textcolor{CTrap}{易错八：参数题只管增减，不管能否取值}}

比如 $\ln(x^2-ax+1)$ 在 $(2,+\infty)$ 上递增，
除了内层二次函数在这段递增，还必须保证每个 $x>2$ 时真数都大于 $0$。
漏掉这一点可能错误地得到 $a\le4$；正确结果是 $a\le\frac52$。

\medskip
\textbf{\textcolor{CTrap}{考前逐项自查}}

\textbf{有定义吗？} \quad
\textbf{内层同向吗？} \quad
\textbf{外层在这些输入上同向吗？} \quad
\textbf{区间写回 $x$ 了吗？}
\end{trapbox}
